% GENERATED FILE. Edit canonical lessons, then run npm run generate:books.
% canonical-source-hash: fnv1a64:789a523061f98394
% canonical-lessons: PA-C228-nichorna, PA-C228-sukauna, PA-C228-tangna, PA-C228-rassi, PA-C228-punjhna

\chapter{Drying the Clothes}
\label{ch:pa-a2-drying-the-clothes}
\hlchaptermodality{228}

\begin{chapteropening}
\textbf{By the end of this chapter:} \emph{I can say how I wring and dry the washing.}

Complete the last lesson of chapter 228: I can say how I wring and dry the washing.
\end{chapteropening}

\section[\texorpdfstring{\pa{ਨਿਚੋੜਨਾ} (nichoṛnā)}{ਨਿਚੋੜਨਾ (nichoṛnā)}]{\pa{ਨਿਚੋੜਨਾ} (nichoṛnā) — to wring}

\label{lesson:PA-C228-nichorna}



\begin{quote}
Before the new one: say the Punjabi for grime, then the Punjabi for a tap.
\end{quote}

\subsection*{You'll want to know: \pa{ਨਿਚੋੜਨਾ}}
\textbf{\pa{ਨਿਚੋੜਨਾ}} — \emph{nichoṛnā} — \textquotedblleft{}to wring\textquotedblright{}.

\begin{quote}
\pa{ਮੈਂ} \pa{ਗਿੱਲੇ} \pa{ਕੱਪੜੇ} \pa{ਨਿਚੋੜਦੀ} \pa{ਹਾਂ।}
\end{quote}

\begin{quote}
\emph{maiṁ gille kappṛe nichoṛdī hāṁ.}
\end{quote}

\begin{quote}
I wring out the wet clothes.
\end{quote}

Past (I …): \pa{ਮੈਂ} \pa{ਤੌਲੀਆ} \pa{ਨਿਚੋੜਿਆ}, \emph{maiṁ taulīā nichoṛiā}.

\begin{cousinweb}
\textbf{\pa{ਨਿਚੋੜਨਾ}} (\emph{nichoṛnā}) is inherited from a reconstructed \emph{*niścōṭayati}, \textquotedblleft{}presses out\textquotedblright{}; CDIAL lists Punjabi \emph{nicoṛṇā}, \textquotedblleft{}to wring\textquotedblright{} (CDIAL 7449).
\end{cousinweb}

\begin{grammarlens}[title={What you've built}]
One more word for saying what you do.
\end{grammarlens}

\subsection*{Your turn}
\begin{itemize}
\raggedright
  \item \emph{Say it:} \emph{nichoṛnā}
  \item \emph{Say it:} \emph{nichoṛnā}, once more
  \item \emph{Say it:} \emph{maiṁ gille kappṛe nichoṛdī hāṁ.} — I wring out the wet clothes.
  \item \emph{Recall:} say the Punjabi for grime, then the Punjabi for a tap, then say \emph{nichoṛnā} again
\end{itemize}

\begin{tcolorbox}[breakable,colback=teal!4,colframe=teal!35!black,title={Before you move on}]
What does \pa{ਨਿਚੋੜਨਾ} mean? (\textquotedblleft{}To wring\textquotedblright{}.) Say it once more.
\end{tcolorbox}
\section[\texorpdfstring{\pa{ਸੁਕਾਉਣਾ} (sukāuṇā)}{ਸੁਕਾਉਣਾ (sukāuṇā)}]{\pa{ਸੁਕਾਉਣਾ} (sukāuṇā) — to dry something}

\label{lesson:PA-C228-sukauna}



\begin{quote}
Before the new one: say the Punjabi for a tap, then the Punjabi for to wring.
\end{quote}

\subsection*{You'll want to know: \pa{ਸੁਕਾਉਣਾ}}
\textbf{\pa{ਸੁਕਾਉਣਾ}} — \emph{sukāuṇā} — \textquotedblleft{}to dry something\textquotedblright{}.

\begin{quote}
\pa{ਅਸੀਂ} \pa{ਕੱਪੜੇ} \pa{ਛੱਤ} \pa{ਉੱਤੇ} \pa{ਸੁਕਾਉਂਦੇ} \pa{ਹਾਂ।}
\end{quote}

\begin{quote}
\emph{asīṁ kappṛe chhatt utte sukāunde hāṁ.}
\end{quote}

\begin{quote}
We dry the clothes on the roof.
\end{quote}

Past (I …): \pa{ਮੈਂ} \pa{ਕੱਪੜੇ} \pa{ਸੁਕਾਏ}, \emph{maiṁ kappṛe sukāe}.

\begin{cousinweb}
\textbf{\pa{ਸੁਕਾਉਣਾ}} (\emph{sukāuṇā}) is the causative of a verb inherited from a reconstructed \emph{*śuṣkati}, \textquotedblleft{}becomes dry\textquotedblright{}; CDIAL lists Punjabi \emph{sukkṇā}, \textquotedblleft{}dries up\textquotedblright{}, and its causative \emph{sukhāuṇā} (CDIAL 12552).
\end{cousinweb}

\begin{grammarlens}[title={What you've built}]
One more word for saying what you do.
\end{grammarlens}

\subsection*{Your turn}
\begin{itemize}
\raggedright
  \item \emph{Say it:} \emph{sukāuṇā}
  \item \emph{Say it:} \emph{sukāuṇā}, once more
  \item \emph{Say it:} \emph{asīṁ kappṛe chhatt utte sukāunde hāṁ.} — We dry the clothes on the roof.
  \item \emph{Recall:} say the Punjabi for a tap, then the Punjabi for to wring, then say \emph{sukāuṇā} again
  \item \emph{Write it:} copy \textbf{\pa{ਸੁਕਾਉਣਾ}} once, checking the whole word
\end{itemize}

\begin{tcolorbox}[breakable,colback=teal!4,colframe=teal!35!black,title={Before you move on}]
What does \pa{ਸੁਕਾਉਣਾ} mean? (\textquotedblleft{}To dry something\textquotedblright{}.) Say it once more.
\end{tcolorbox}
\section[\texorpdfstring{\pa{ਟੰਗਣਾ} (ṭaṅgṇā)}{ਟੰਗਣਾ (ṭaṅgṇā)}]{\pa{ਟੰਗਣਾ} (ṭaṅgṇā) — to hang up}

\label{lesson:PA-C228-tangna}



\begin{quote}
Before the new one: say the Punjabi for to wring, then the Punjabi for to dry something.
\end{quote}

\subsection*{You'll want to know: \pa{ਟੰਗਣਾ}}
\textbf{\pa{ਟੰਗਣਾ}} — \emph{ṭaṅgṇā} — \textquotedblleft{}to hang up\textquotedblright{}.

\begin{quote}
\pa{ਮੈਂ} \pa{ਕਮੀਜ਼} \pa{ਕਿੱਲ} \pa{ਉੱਤੇ} \pa{ਟੰਗਦਾ} \pa{ਹਾਂ।}
\end{quote}

\begin{quote}
\emph{maiṁ kamīz kill utte ṭaṅgdā hāṁ.}
\end{quote}

\begin{quote}
I hang my shirt on a nail.
\end{quote}

Past (I …): \pa{ਮੈਂ} \pa{ਕਮੀਜ਼} \pa{ਟੰਗੀ}, \emph{maiṁ kamīz ṭaṅgī}.

\begin{cousinweb}
\textbf{\pa{ਟੰਗਣਾ}} (\emph{ṭaṅgṇā}) is inherited from a reconstructed \emph{*ṭaṅg}, \textquotedblleft{}hang up\textquotedblright{}; CDIAL lists Punjabi \emph{ṭaṅgṇā}, \textquotedblleft{}to hang up\textquotedblright{} (CDIAL 5436).
\end{cousinweb}

\begin{grammarlens}[title={What you've built}]
One more word for saying what you do.
\end{grammarlens}

\subsection*{Your turn}
\begin{itemize}
\raggedright
  \item \emph{Say it:} \emph{ṭaṅgṇā}
  \item \emph{Say it:} \emph{ṭaṅgṇā}, once more
  \item \emph{Say it:} \emph{maiṁ kamīz kill utte ṭaṅgdā hāṁ.} — I hang my shirt on a nail.
  \item \emph{Recall:} say the Punjabi for to wring, then the Punjabi for to dry something, then say \emph{ṭaṅgṇā} again
  \item \emph{Write it:} copy \textbf{\pa{ਟੰਗਣਾ}} once, checking the whole word
\end{itemize}

\begin{tcolorbox}[breakable,colback=teal!4,colframe=teal!35!black,title={Before you move on}]
What does \pa{ਟੰਗਣਾ} mean? (\textquotedblleft{}To hang up\textquotedblright{}.) Say it once more.
\end{tcolorbox}
\section[\texorpdfstring{\pa{ਰੱਸੀ} (rassī)}{ਰੱਸੀ (rassī)}]{\pa{ਰੱਸੀ} (rassī) — a rope; a washing line}

\label{lesson:PA-C228-rassi}



\begin{quote}
Before the new one: say the Punjabi for to dry something, then the Punjabi for to hang up.
\end{quote}

\subsection*{You'll want to know: \pa{ਰੱਸੀ}}
\textbf{\pa{ਰੱਸੀ}} — \emph{rassī} — \textquotedblleft{}a rope; a washing line\textquotedblright{}.

\begin{quote}
\pa{ਗਿੱਲੇ} \pa{ਕੱਪੜੇ} \pa{ਰੱਸੀ} \pa{ਉੱਤੇ} \pa{ਹਨ।}
\end{quote}

\begin{quote}
\emph{gille kappṛe rassī utte han.}
\end{quote}

\begin{quote}
The wet clothes are on the line.
\end{quote}

A feminine noun.

\begin{cousinweb}
\textbf{\pa{ਰੱਸੀ}} (\emph{rassī}) is inherited from Sanskrit \emph{raśmi}, \textquotedblleft{}rope, bridle, ray of light\textquotedblright{}; CDIAL lists Punjabi \emph{rassī}, \textquotedblleft{}rope\textquotedblright{} (CDIAL 10648).
\end{cousinweb}

\begin{grammarlens}[title={What you've built}]
One more word for saying what you do.
\end{grammarlens}

\subsection*{Your turn}
\begin{itemize}
\raggedright
  \item \emph{Say it:} \emph{rassī}
  \item \emph{Say it:} \emph{rassī}, once more
  \item \emph{Say it:} \emph{gille kappṛe rassī utte han.} — The wet clothes are on the line.
  \item \emph{Recall:} say the Punjabi for to dry something, then the Punjabi for to hang up, then say \emph{rassī} again
  \item \emph{Write it:} copy \textbf{\pa{ਰੱਸੀ}} once, checking the whole word
\end{itemize}

\begin{tcolorbox}[breakable,colback=teal!4,colframe=teal!35!black,title={Before you move on}]
What does \pa{ਰੱਸੀ} mean? (\textquotedblleft{}A rope; a washing line\textquotedblright{}.) Say it once more.
\end{tcolorbox}
\section[\texorpdfstring{\pa{ਪੂੰਝਣਾ} (pūnjhṇā)}{ਪੂੰਝਣਾ (pūnjhṇā)}]{\pa{ਪੂੰਝਣਾ} (pūnjhṇā) — to wipe}

\label{lesson:PA-C228-punjhna}



\begin{quote}
Before the new one: say the Punjabi for to hang up, then the Punjabi for a rope.
\end{quote}

\subsection*{You'll want to know: \pa{ਪੂੰਝਣਾ}}
\textbf{\pa{ਪੂੰਝਣਾ}} — \emph{pūnjhṇā} — \textquotedblleft{}to wipe\textquotedblright{}.

\begin{quote}
\pa{ਮੈਂ} \pa{ਤੌਲੀਏ} \pa{ਨਾਲ} \pa{ਮੂੰਹ} \pa{ਪੂੰਝਦਾ} \pa{ਹਾਂ।}
\end{quote}

\begin{quote}
\emph{maiṁ taulīe nāl mūnh pūnjhdā hāṁ.}
\end{quote}

\begin{quote}
I wipe my face with a towel.
\end{quote}

Past (I …): \pa{ਮੈਂ} \pa{ਮੇਜ਼} \pa{ਪੂੰਝਿਆ}, \emph{maiṁ mez pūnjhiā}.

\begin{cousinweb}
\textbf{\pa{ਪੂੰਝਣਾ}} (\emph{pūnjhṇā}) is inherited from Sanskrit \emph{prōñchati}, \textquotedblleft{}wipes off\textquotedblright{}; CDIAL lists Punjabi \emph{pū̃jhṇā}, \textquotedblleft{}wipes off\textquotedblright{} (CDIAL 9011).
\end{cousinweb}

\begin{grammarlens}[title={What you've built}]
One more word for saying what you do.
\end{grammarlens}

\subsection*{Your turn}
\begin{itemize}
\raggedright
  \item \emph{Say it:} \emph{pūnjhṇā}
  \item \emph{Say it:} \emph{pūnjhṇā}, once more
  \item \emph{Say it:} \emph{maiṁ taulīe nāl mūnh pūnjhdā hāṁ.} — I wipe my face with a towel.
  \item \emph{Recall:} say the Punjabi for to hang up, then the Punjabi for a rope, then say \emph{pūnjhṇā} again
  \item \emph{Write it:} copy \textbf{\pa{ਪੂੰਝਣਾ}} once, checking the whole word
\end{itemize}

\begin{tcolorbox}[breakable,colback=teal!4,colframe=teal!35!black,title={Before you move on}]
What does \pa{ਪੂੰਝਣਾ} mean? (\textquotedblleft{}To wipe\textquotedblright{}.) Say it once more.
\end{tcolorbox}
